\documentclass[11pt,titlepage]{article}

\usepackage{kotex}
\usepackage{amsmath}
\usepackage{booktabs}
\usepackage{geometry}
\usepackage{hyperref}
\usepackage{cite}
\usepackage{tabularx}
\usepackage{array}
\usepackage{longtable}

\geometry{a4paper, margin=2.5cm}
\hypersetup{hidelinks,pdftitle={Paper UI Color System},pdfauthor={Jaewon Baek}}
\setlength{\parskip}{0.35em}
\setlength{\parindent}{1em}
\renewcommand{\arraystretch}{1.18}
\newcolumntype{Y}{>{\raggedright\arraybackslash}X}

\title{Paper UI Color System\\[0.4em]
\large OKLCH 팔레트에서 실제 화면까지}
\author{Jaewon Baek}
\date{2026년 9월 25일}

\begin{document}

\hypersetup{pageanchor=false}
\maketitle
\hypersetup{pageanchor=true}
\pagenumbering{arabic}

\newpage
\begin{abstract}
Paper UI의 색상 시스템은 개별 컴포넌트에 색을 지정하는 방식이 아니라, OKLCH 원천 팔레트에서 생성한 HEX palette token을 semantic token으로 변환하고 최종 UI가 semantic CSS variable을 사용하게 하는 단방향 구조다. 이 보고서는 현재 구현을 재현하고 검증할 수 있도록 팔레트 전체 좌표와 HEX 결과, 생성 알고리즘, semantic의 정확한 step 선택, 모든 accent variant의 배경·전경·경계·hover recipe, interaction alpha, CSS 전달 경로, 컴포넌트 상태, 테스트 계약과 알려진 예외를 기록한다. 구조적 정합성은 상당 부분 확인되지만 실제 배경·글자 조합의 명도 대비를 모두 평가하는 자동 검증은 별도의 보강 과제로 남는다.
\end{abstract}

\newpage
\renewcommand{\contentsname}{목차}
\tableofcontents
\newpage

\section{서론}

디자인 시스템에서 색은 팔레트에 예쁜 색을 모으는 일만으로 완성되지 않는다. 같은 회색이 한 화면에서는 페이지 바탕이고 다른 곳에서는 글자나 경계선이 될 수 있으며, 라이트와 다크 테마에서는 그 관계가 다시 달라진다. 역할을 컴포넌트마다 직접 정하면 색의 의미가 분산되고, 한 색을 조정할 때 수정 위치도 불어난다.

Paper UI는 이 문제를 네 단계의 단방향 계약으로 푼다. 작성자가 관리하는 OKLCH source를 HEX palette token으로 변환하고, semantic token은 palette token만 선택한다. 마지막으로 CSS와 앱은 semantic token만 참조한다. 따라서 색의 출처와 화면에서의 의미가 분리되며, 팔레트 조정은 한 곳에서 하고 사용처는 역할 이름을 통해 유지할 수 있다\cite{architecture,palette,ladder,semantics}.

\section{설계 원칙과 색상 흐름}

Paper UI의 색은 화면의 구조와 의미를 나누어 맡는다. neutral은 surface, border, ink의 명도 위계를 만들고, blue·green·amber·red는 정보·성공·경고·오류처럼 의미가 필요한 위치에서만 사용한다. 주요 동작의 기본색도 임의의 브랜드 색이 아니라 neutral 계열에 둔다. 넓은 면을 accent로 장식하기보다 중립색의 밝기와 경계로 밀도를 만들고, accent는 작은 상태 표식과 상호작용에 집중한다\cite{tokensdoc}.

구현은 아래의 네 단계로 제한된다.

\begin{center}
\texttt{OKLCH source $\rightarrow$ HEX palette tokens $\rightarrow$ semantic systems $\rightarrow$ CSS variables / UI}
\end{center}

각 층은 다음 층만을 위한 입력을 제공한다. Palette source에는 화면 역할이나 CSS 이름이 없고, 생성기는 semantic 판단을 하지 않는다. Semantic leaf는 다른 semantic leaf를 참조하지 않으며, component·resolver·앱은 source 좌표나 palette generator를 직접 읽지 않는다. 이 경계를 지키면 색 하나의 변경이 어떤 화면 요소를 바꾸는지 추적할 수 있다\cite{architecture}.

\section{OKLCH source와 HEX palette token}

\subsection{Palette source}

\texttt{palette.ts}는 라이트와 다크 테마별 OKLCH anchor와 interaction alpha 설정의 단일 작성 지점이다. Neutral ramp는 13단계 \texttt{0, 50, 100, 200, 300, 400, 500, 600, 700, 800, 900, 925, 950}로, accent ramp는 11단계 \texttt{50--950}로 구성된다. 네 accent family 모두 50, 100, 200, 400, 600, 700, 950 anchor를 갖고, 빠진 단계는 생성기에서 보간한다\cite{palette}.

Neutral step 번호는 테마와 무관하게 밝은 쪽에서 어두운 쪽으로 진행한다. 따라서 \texttt{neutral.0}은 white에 가까운 끝점이고 \texttt{neutral.950}은 near-black 끝점이다. 이번 변경은 기존에 검토한 두 ramp의 색상 자체를 바꾸지 않고 주소만 이 방향으로 재색인했다. 라이트 neutral은 hue 286°와 낮은 chroma 0.003의 cool-neutral이며, 다크 neutral은 chroma 0의 charcoal이다. 기준 anchor 사이 단계는 순서상 같은 간격으로 OKLCH lightness를 선형 보간한다\cite{palette}.

Accent hue는 테마별 anchor로 보존한다. 예를 들어 라이트 blue는 약 254°를 중심으로, 다크 blue는 약 258°를 중심으로 한다. 다크 amber와 red의 hue도 라이트와 같게 고정하지 않는다. 이는 양 테마가 동일한 수치 좌표를 반복하는 대신, 실제 화면에서 각 family의 정체성이 유지되도록 anchor를 따로 검토할 수 있게 한다\cite{palette}.

\subsection{모든 generated palette 값}
Neutral 및 accent의 모든 step은 L0에서 OKLCH로 완성된 좌표를 갖고, L1 generator가 동일 좌표를 HEX로 바꾼다. 아래 표의 좌표는 보간까지 끝난 실제 입력값이며, HEX는 gamut mapping과 8-bit sRGB 양자화를 거친 출력이다. 소수 넷째 자리까지 표기했으므로 HEX 생성 시 사용되는 원래 좌표의 마지막 자리는 반올림되어 보이지 않을 수 있다. Neutral 번호는 light-to-dark 순서다. Accent 번호는 neutral처럼 공통 밝기 순서가 아니다. light accent는 낮은 step에서 옅은 wash, dark accent는 낮은 step에서 어두운 wash를 시작하며, 두 테마 모두 semantic slot이 같은 step family를 고른다.

\small
\setlength{\tabcolsep}{3pt}
\begin{longtable}{@{}r r r r l r r r l@{}}
\caption{Neutral 전체 ramp의 source OKLCH와 생성 HEX(light/dark)}\label{tab:full-neutral}\\
\toprule
Step & \multicolumn{4}{c}{Light: L C H HEX} & \multicolumn{4}{c}{Dark: L C H HEX}\\
 & L & C & H & HEX & L & C & H & HEX\\
\midrule
\endfirsthead
\multicolumn{9}{l}{\tablename\ \thetable\ (continued)}\\
\toprule
Step & \multicolumn{4}{c}{Light: L C H HEX} & \multicolumn{4}{c}{Dark: L C H HEX}\\
 & L & C & H & HEX & L & C & H & HEX\\
\midrule
\endhead
\midrule
\multicolumn{9}{r}{Continued on next page}\\
\endfoot
\bottomrule
\endlastfoot
0 & \texttt{1.0000} & \texttt{0.0030} & \texttt{286.0} & \texttt{\#ffffff} & \texttt{0.9702} & \texttt{0.0000} & \texttt{0.0} & \texttt{\#f5f5f5} \\
50 & \texttt{0.9912} & \texttt{0.0030} & \texttt{286.0} & \texttt{\#fcfcfe} & \texttt{0.9310} & \texttt{0.0000} & \texttt{0.0} & \texttt{\#e8e8e8} \\
100 & \texttt{0.9824} & \texttt{0.0030} & \texttt{286.0} & \texttt{\#f9f9fb} & \texttt{0.8576} & \texttt{0.0000} & \texttt{0.0} & \texttt{\#d0d0d0} \\
200 & \texttt{0.9485} & \texttt{0.0030} & \texttt{286.0} & \texttt{\#eeeef0} & \texttt{0.7476} & \texttt{0.0000} & \texttt{0.0} & \texttt{\#adadad} \\
300 & \texttt{0.9145} & \texttt{0.0030} & \texttt{286.0} & \texttt{\#e2e2e5} & \texttt{0.6467} & \texttt{0.0000} & \texttt{0.0} & \texttt{\#8e8e8e} \\
400 & \texttt{0.8806} & \texttt{0.0030} & \texttt{286.0} & \texttt{\#d7d7d9} & \texttt{0.5555} & \texttt{0.0000} & \texttt{0.0} & \texttt{\#737373} \\
500 & \texttt{0.8467} & \texttt{0.0030} & \texttt{286.0} & \texttt{\#ccccce} & \texttt{0.4748} & \texttt{0.0000} & \texttt{0.0} & \texttt{\#5c5c5c} \\
600 & \texttt{0.7616} & \texttt{0.0030} & \texttt{286.0} & \texttt{\#b1b1b3} & \texttt{0.3867} & \texttt{0.0000} & \texttt{0.0} & \texttt{\#444444} \\
700 & \texttt{0.6765} & \texttt{0.0030} & \texttt{286.0} & \texttt{\#979799} & \texttt{0.3250} & \texttt{0.0000} & \texttt{0.0} & \texttt{\#343434} \\
800 & \texttt{0.5914} & \texttt{0.0030} & \texttt{286.0} & \texttt{\#7d7d7f} & \texttt{0.2891} & \texttt{0.0000} & \texttt{0.0} & \texttt{\#2b2b2b} \\
900 & \texttt{0.5063} & \texttt{0.0030} & \texttt{286.0} & \texttt{\#656567} & \texttt{0.2478} & \texttt{0.0000} & \texttt{0.0} & \texttt{\#212121} \\
925 & \texttt{0.3885} & \texttt{0.0030} & \texttt{286.0} & \texttt{\#444446} & \texttt{0.2046} & \texttt{0.0000} & \texttt{0.0} & \texttt{\#171717} \\
950 & \texttt{0.2707} & \texttt{0.0030} & \texttt{286.0} & \texttt{\#262628} & \texttt{0.1730} & \texttt{0.0000} & \texttt{0.0} & \texttt{\#101010} \\
\end{longtable}
\begin{longtable}{@{}r r r r l r r r l@{}}
\caption{Blue accent 전체 ramp의 source OKLCH와 생성 HEX(light/dark)}\label{tab:accent-blue}\\
\toprule
Step & \multicolumn{4}{c}{Light: L C H HEX} & \multicolumn{4}{c}{Dark: L C H HEX}\\
 & L & C & H & HEX & L & C & H & HEX\\
\midrule
\endfirsthead
\multicolumn{9}{l}{\tablename\ \thetable\ (continued)}\\
\toprule
Step & \multicolumn{4}{c}{Light: L C H HEX} & \multicolumn{4}{c}{Dark: L C H HEX}\\
 & L & C & H & HEX & L & C & H & HEX\\
\midrule
\endhead
\midrule
\multicolumn{9}{r}{Continued on next page}\\
\endfoot
\bottomrule
\endlastfoot
50 & \texttt{0.9805} & \texttt{0.0079} & \texttt{253.9} & \texttt{\#f5f9fe} & \texttt{0.2261} & \texttt{0.0478} & \texttt{257.9} & \texttt{\#0d1c32} \\
100 & \texttt{0.9304} & \texttt{0.0263} & \texttt{254.0} & \texttt{\#dce9fa} & \texttt{0.2561} & \texttt{0.0708} & \texttt{258.2} & \texttt{\#0b2244} \\
200 & \texttt{0.9039} & \texttt{0.0324} & \texttt{254.0} & \texttt{\#d1e1f5} & \texttt{0.3301} & \texttt{0.0784} & \texttt{258.2} & \texttt{\#1a355d} \\
300 & \texttt{0.7837} & \texttt{0.0875} & \texttt{255.4} & \texttt{\#93bbf0} & \texttt{0.4343} & \texttt{0.0849} & \texttt{258.1} & \texttt{\#32517f} \\
400 & \texttt{0.6636} & \texttt{0.1426} & \texttt{256.7} & \texttt{\#5694e9} & \texttt{0.5384} & \texttt{0.0915} & \texttt{258.1} & \texttt{\#4c6fa3} \\
500 & \texttt{0.6048} & \texttt{0.1789} & \texttt{259.8} & \texttt{\#3a7dea} & \texttt{0.5993} & \texttt{0.1204} & \texttt{257.9} & \texttt{\#5080c7} \\
600 & \texttt{0.5461} & \texttt{0.2152} & \texttt{262.9} & \texttt{\#2563eb} & \texttt{0.6601} & \texttt{0.1493} & \texttt{257.8} & \texttt{\#5492ec} \\
700 & \texttt{0.4882} & \texttt{0.2172} & \texttt{264.4} & \texttt{\#1d4ed8} & \texttt{0.7882} & \texttt{0.1080} & \texttt{257.7} & \texttt{\#8ebcff} \\
800 & \texttt{0.4206} & \texttt{0.1797} & \texttt{264.1} & \texttt{\#1740ad} & \texttt{0.8472} & \texttt{0.0775} & \texttt{258.3} & \texttt{\#afcfff} \\
900 & \texttt{0.3529} & \texttt{0.1423} & \texttt{263.9} & \texttt{\#123284} & \texttt{0.9061} & \texttt{0.0470} & \texttt{258.8} & \texttt{\#cee1ff} \\
950 & \texttt{0.2853} & \texttt{0.1048} & \texttt{263.6} & \texttt{\#0d255d} & \texttt{0.9651} & \texttt{0.0165} & \texttt{259.4} & \texttt{\#edf4ff} \\
\end{longtable}

\begin{longtable}{@{}r r r r l r r r l@{}}
\caption{Green accent 전체 ramp의 source OKLCH와 생성 HEX(light/dark)}\label{tab:accent-green}\\
\toprule
Step & \multicolumn{4}{c}{Light: L C H HEX} & \multicolumn{4}{c}{Dark: L C H HEX}\\
 & L & C & H & HEX & L & C & H & HEX\\
\midrule
\endfirsthead
\multicolumn{9}{l}{\tablename\ \thetable\ (continued)}\\
\toprule
Step & \multicolumn{4}{c}{Light: L C H HEX} & \multicolumn{4}{c}{Dark: L C H HEX}\\
 & L & C & H & HEX & L & C & H & HEX\\
\midrule
\endhead
\midrule
\multicolumn{9}{r}{Continued on next page}\\
\endfoot
\bottomrule
\endlastfoot
50 & \texttt{0.9867} & \texttt{0.0107} & \texttt{158.9} & \texttt{\#f5fdf8} & \texttt{0.2250} & \texttt{0.0475} & \texttt{153.7} & \texttt{\#062211} \\
100 & \texttt{0.9319} & \texttt{0.0355} & \texttt{158.8} & \texttt{\#d6f0e0} & \texttt{0.2558} & \texttt{0.0700} & \texttt{153.2} & \texttt{\#002c13} \\
200 & \texttt{0.9259} & \texttt{0.0437} & \texttt{158.8} & \texttt{\#cff0dc} & \texttt{0.3285} & \texttt{0.0774} & \texttt{153.2} & \texttt{\#094021} \\
300 & \texttt{0.7865} & \texttt{0.0955} & \texttt{155.3} & \texttt{\#86cc9f} & \texttt{0.4284} & \texttt{0.0854} & \texttt{153.6} & \texttt{\#225d39} \\
400 & \texttt{0.6472} & \texttt{0.1473} & \texttt{151.8} & \texttt{\#35a75f} & \texttt{0.5284} & \texttt{0.0934} & \texttt{154.0} & \texttt{\#3a7b52} \\
500 & \texttt{0.5873} & \texttt{0.1422} & \texttt{150.9} & \texttt{\#26934e} & \texttt{0.5945} & \texttt{0.1220} & \texttt{153.5} & \texttt{\#38935b} \\
600 & \texttt{0.5273} & \texttt{0.1371} & \texttt{150.1} & \texttt{\#15803d} & \texttt{0.6606} & \texttt{0.1505} & \texttt{152.9} & \texttt{\#32ac64} \\
700 & \texttt{0.4479} & \texttt{0.1083} & \texttt{151.3} & \texttt{\#166534} & \texttt{0.7894} & \texttt{0.1094} & \texttt{153.1} & \texttt{\#81cf99} \\
800 & \texttt{0.3842} & \texttt{0.0913} & \texttt{152.3} & \texttt{\#0f512a} & \texttt{0.8489} & \texttt{0.0817} & \texttt{152.9} & \texttt{\#a5ddb4} \\
900 & \texttt{0.3204} & \texttt{0.0742} & \texttt{153.2} & \texttt{\#0a3d20} & \texttt{0.9085} & \texttt{0.0541} & \texttt{152.6} & \texttt{\#c7ecd0} \\
950 & \texttt{0.2567} & \texttt{0.0572} & \texttt{154.1} & \texttt{\#062b16} & \texttt{0.9680} & \texttt{0.0264} & \texttt{152.4} & \texttt{\#e8faec} \\
\end{longtable}

\begin{longtable}{@{}r r r r l r r r l@{}}
\caption{Amber accent 전체 ramp의 source OKLCH와 생성 HEX(light/dark)}\label{tab:accent-amber}\\
\toprule
Step & \multicolumn{4}{c}{Light: L C H HEX} & \multicolumn{4}{c}{Dark: L C H HEX}\\
 & L & C & H & HEX & L & C & H & HEX\\
\midrule
\endfirsthead
\multicolumn{9}{l}{\tablename\ \thetable\ (continued)}\\
\toprule
Step & \multicolumn{4}{c}{Light: L C H HEX} & \multicolumn{4}{c}{Dark: L C H HEX}\\
 & L & C & H & HEX & L & C & H & HEX\\
\midrule
\endhead
\midrule
\multicolumn{9}{r}{Continued on next page}\\
\endfoot
\bottomrule
\endlastfoot
50 & \texttt{0.9882} & \texttt{0.0110} & \texttt{89.7} & \texttt{\#fefbf3} & \texttt{0.2261} & \texttt{0.0479} & \texttt{74.2} & \texttt{\#291800} \\
100 & \texttt{0.9323} & \texttt{0.0524} & \texttt{95.2} & \texttt{\#f3e9c2} & \texttt{0.2561} & \texttt{0.0654} & \texttt{70.9} & \texttt{\#331e00} \\
200 & \texttt{0.9262} & \texttt{0.0662} & \texttt{95.2} & \texttt{\#f4e7b5} & \texttt{0.3304} & \texttt{0.0712} & \texttt{70.9} & \texttt{\#4c2e00} \\
300 & \texttt{0.7970} & \texttt{0.1003} & \texttt{84.5} & \texttt{\#dbb86f} & \texttt{0.4363} & \texttt{0.0809} & \texttt{72.9} & \texttt{\#6c4a16} \\
400 & \texttt{0.6678} & \texttt{0.1343} & \texttt{73.8} & \texttt{\#c5861b} & \texttt{0.5422} & \texttt{0.0905} & \texttt{74.9} & \texttt{\#8e672c} \\
500 & \texttt{0.6116} & \texttt{0.1399} & \texttt{61.4} & \texttt{\#be6d06} & \texttt{0.6018} & \texttt{0.1154} & \texttt{74.3} & \texttt{\#a97520} \\
600 & \texttt{0.5553} & \texttt{0.1455} & \texttt{49.0} & \texttt{\#b45309} & \texttt{0.6613} & \texttt{0.1404} & \texttt{73.7} & \texttt{\#c58300} \\
700 & \texttt{0.4732} & \texttt{0.1247} & \texttt{46.2} & \texttt{\#92400e} & \texttt{0.7896} & \texttt{0.1100} & \texttt{74.8} & \texttt{\#e4b066} \\
800 & \texttt{0.4101} & \texttt{0.1045} & \texttt{47.9} & \texttt{\#76350a} & \texttt{0.8477} & \texttt{0.0868} & \texttt{78.8} & \texttt{\#ecc78c} \\
900 & \texttt{0.3469} & \texttt{0.0843} & \texttt{49.6} & \texttt{\#5b2a08} & \texttt{0.9058} & \texttt{0.0636} & \texttt{82.7} & \texttt{\#f5ddb0} \\
950 & \texttt{0.2838} & \texttt{0.0641} & \texttt{51.3} & \texttt{\#421f07} & \texttt{0.9639} & \texttt{0.0404} & \texttt{86.7} & \texttt{\#fff2d5} \\
\end{longtable}

\begin{longtable}{@{}r r r r l r r r l@{}}
\caption{Red accent 전체 ramp의 source OKLCH와 생성 HEX(light/dark)}\label{tab:accent-red}\\
\toprule
Step & \multicolumn{4}{c}{Light: L C H HEX} & \multicolumn{4}{c}{Dark: L C H HEX}\\
 & L & C & H & HEX & L & C & H & HEX\\
\midrule
\endfirsthead
\multicolumn{9}{l}{\tablename\ \thetable\ (continued)}\\
\toprule
Step & \multicolumn{4}{c}{Light: L C H HEX} & \multicolumn{4}{c}{Dark: L C H HEX}\\
 & L & C & H & HEX & L & C & H & HEX\\
\midrule
\endhead
\midrule
\multicolumn{9}{r}{Continued on next page}\\
\endfoot
\bottomrule
\endlastfoot
50 & \texttt{0.9808} & \texttt{0.0064} & \texttt{17.3} & \texttt{\#fdf7f7} & \texttt{0.2265} & \texttt{0.0478} & \texttt{23.6} & \texttt{\#2f1211} \\
100 & \texttt{0.9304} & \texttt{0.0250} & \texttt{17.7} & \texttt{\#f8e2e2} & \texttt{0.2562} & \texttt{0.0711} & \texttt{21.3} & \texttt{\#3f1113} \\
200 & \texttt{0.9053} & \texttt{0.0312} & \texttt{17.7} & \texttt{\#f4d8d8} & \texttt{0.3300} & \texttt{0.0788} & \texttt{21.3} & \texttt{\#572223} \\
300 & \texttt{0.7939} & \texttt{0.0975} & \texttt{20.1} & \texttt{\#f4a3a2} & \texttt{0.4397} & \texttt{0.0862} & \texttt{21.3} & \texttt{\#7b3d3d} \\
400 & \texttt{0.6825} & \texttt{0.1639} & \texttt{22.4} & \texttt{\#ed6969} & \texttt{0.5495} & \texttt{0.0936} & \texttt{21.4} & \texttt{\#a15a59} \\
500 & \texttt{0.6298} & \texttt{0.1895} & \texttt{24.9} & \texttt{\#e54c4b} & \texttt{0.6045} & \texttt{0.1220} & \texttt{21.5} & \texttt{\#c06161} \\
600 & \texttt{0.5771} & \texttt{0.2152} & \texttt{27.3} & \texttt{\#dc2626} & \texttt{0.6595} & \texttt{0.1505} & \texttt{21.6} & \texttt{\#df6768} \\
700 & \texttt{0.5054} & \texttt{0.1905} & \texttt{27.5} & \texttt{\#b91c1c} & \texttt{0.7901} & \texttt{0.1093} & \texttt{21.9} & \texttt{\#f99e9b} \\
800 & \texttt{0.4362} & \texttt{0.1586} & \texttt{26.7} & \texttt{\#951a1b} & \texttt{0.8444} & \texttt{0.0806} & \texttt{23.7} & \texttt{\#fdb8b3} \\
900 & \texttt{0.3671} & \texttt{0.1266} & \texttt{25.8} & \texttt{\#731718} & \texttt{0.8987} & \texttt{0.0520} & \texttt{25.5} & \texttt{\#ffd1cc} \\
950 & \texttt{0.2979} & \texttt{0.0947} & \texttt{25.0} & \texttt{\#531313} & \texttt{0.9530} & \texttt{0.0233} & \texttt{27.3} & \texttt{\#ffeae7} \\
\end{longtable}

\setlength{\tabcolsep}{6pt}
\normalsize

\subsection{HEX 생성과 gamut mapping}

\texttt{color-ladder.ts}는 OKLCH를 선형 sRGB로 변환하고, sRGB 범위 안에 있으면 해당 색을 사용한다. 범위를 벗어나면 lightness와 hue를 고정한 채 chroma를 이진 탐색으로 줄여 gamut 안에 들어오는 최대 chroma를 찾는다. 최종 채널은 sRGB transfer function을 적용해 8-bit 정수로 양자화하고 6자리 HEX로 출력한다. Accent 중간 단계는 인접한 OKLCH anchor 사이에서 lightness·chroma를 선형 보간하고, hue는 최단 각도로 보간한다\cite{ladder}.

변환은 네 단계로 고정되어 있다. 먼저 hue를 라디안으로 바꾸고 \(a=C\cos h\), \(b=C\sin h\)를 얻는다. 이어 OKLCH의 \((L,a,b)\)에서 선형 LMS의 세제곱근 좌표를 만든다.

\begin{align*}
l' &= L + 0.3963377774a + 0.2158037573b,\\
m' &= L - 0.1055613458a - 0.0638541728b,\\
s' &= L - 0.0894841775a - 1.2914855480b,\\
l &= (l')^3, \quad m=(m')^3, \quad s=(s')^3.
\end{align*}

이 세제곱된 선형 LMS로 선형 sRGB를 계산한다.

\begin{align*}
R_l &= 4.0767416621l - 3.3077115913m + 0.2309699292s,\\
G_l &= -1.2684380046l + 2.6097574011m - 0.3413193965s,\\
B_l &= -0.0041960863l - 0.7034186147m + 1.7076147010s.
\end{align*}

선형 채널이 0.0031308 이하이면 \(12.92x\), 그보다 크면 \(1.055x^{1/2.4}-0.055\)로 sRGB에 인코딩한다. 이후 채널을 0-1 범위에 고정하고 255를 곱해 가장 가까운 정수로 반올림한다. 이 정수 세 개를 2자리 16진수로 이어 붙인 결과가 화면에 전달되는 HEX다. \texttt{\#RRGGBB}는 디스플레이 출력값이고 OKLCH는 그것을 만든 설계 좌표다. 양자화와 gamut mapping 때문에 HEX에서 역산한 좌표가 source와 완전히 같을 필요는 없다.

Gamut mapping은 gamma encoding 전의 선형 RGB 세 채널이 모두 0 이상 1 이하인지 검사한다. 이미 범위 안이면 chroma를 유지한다. 범위 밖이면 \((L,H)\)를 고정하고 chroma 구간 \([0,C_{source}]\)를 24회 이진 탐색해 sRGB 경계 안에 남는 최대 chroma를 선택한다. 매 회차의 중간값을 RGB로 바꿔 세 채널이 범위 안이면 하한을 올리고, 아니면 상한을 내린다. 최종 출력은 마지막 gamut-safe 하한으로 다시 계산한다. 그러므로 extreme lightness의 neutral이나 높은 채도의 accent는 일부 C가 줄어들 수 있지만, tone과 hue는 임의로 이동하지 않는다. 이어 sRGB transfer function을 적용하고 0-1로 clamp한 뒤 8-bit로 양자화한다. Alpha token은 이미 만들어진 RGB HEX에 alpha byte를 덧붙이는 별도 경로다.

Accent 보간은 숫자 자체의 차이가 아니라 \texttt{ACCENT\_STEPS} 배열에서 두 anchor가 차지하는 index 거리로 비율을 정한다. 예를 들어 100과 200 사이의 누락 단계 하나는 anchor 간격을 반으로 나눈다. hue는 최단 각도를 따라 보간되며 0°/360° 경계에서 한 바퀴를 도는 긴 경로를 택하지 않는다. Neutral light ramp는 별도 함수로 L anchor 사이를 배열 순서대로 선형 보간하고 C=0.003, H=286°를 공유한다. Dark neutral은 13단계 좌표를 모두 직접 선언해 C=0의 charcoal ramp를 유지한다.

Interaction과 shadow tint는 palette source의 색에 alpha를 합성한다. 중립색 source step과 opacity로 8자리 HEX palette token을 만들며, status glow는 해당 accent family의 생성된 HEX에 alpha를 붙인다. 이 때문에 semantic layer에 \texttt{rgba()}나 별도 색 계산을 넣지 않아도 투명한 hover, scrim, glow, shadow tint를 생성할 수 있다\cite{palette,ladder}.

Alpha는 source color와 0-1 opacity를 받는다. 출력 alpha byte는 \(\operatorname{round}(255\alpha)\)이며, 이를 RGB HEX 뒤에 붙여 \texttt{\#RRGGBBAA}를 만든다. 따라서 opacity 0.05는 \texttt{0d}, 0.12는 \texttt{1f}, 0.13은 \texttt{21}, 0.14는 \texttt{24}, 0.16은 \texttt{29}, 0.22는 \texttt{38}, 0.32는 \texttt{52}, 0.58은 \texttt{94}로 반올림된다. 코드가 쓰는 CSS 8자리 HEX 표기에서 마지막 두 자리는 RGB가 아니라 alpha다.

Palette step은 색을 식별하는 내부 주소이지, 화면에서의 의미를 직접 뜻하지 않는다. Neutral 숫자는 밝음에서 어두움으로 진행하지만, 의미 계층은 별도의 mapping에서 선택한다. 따라서 \texttt{neutral.100}이 언제나 페이지 바탕이라는 식의 규칙은 없다. 그 역할은 semantic layer가 결정한다.

\section{다섯 semantic color system}

Semantic color는 다섯 system으로 공개된다. Surface는 \texttt{bg}와 \texttt{border}, 나머지는 \texttt{ink}, \texttt{accent}, \texttt{interaction}, \texttt{constants}다. Constants는 shadow tint처럼 palette에서 온 고정 효과색을 보관한다. 각각의 leaf는 generated palette token을 선택할 뿐이며, 다른 semantic branch를 계산 입력으로 삼지 않는다\cite{semantics}.

\subsection{Surface와 ink}

Surface system은 네 배경 단계와 세 경계 단계로 UI의 앞뒤와 구획을 표현한다. Ink system은 같은 neutral family에서 주요·보조·약한 텍스트 색을 고른다. 현재 주요 매핑은 Table~\ref{tab:surface-ink}와 같다. 라이트와 다크는 같은 이름을 사용하되 서로 다른 palette step을 선택하므로 호출부는 테마 분기를 직접 구현하지 않는다\cite{semantics}.

\begin{table}[htbp]
\centering
\caption{Surface 및 ink semantic mapping}
\label{tab:surface-ink}
\small
\begin{tabularx}{\textwidth}{@{}l X X Y@{}}
\toprule
Semantic token & Light step / HEX / OKLCH (L, C, H) & Dark step / HEX / OKLCH (L, C, H) & 역할 \\
\midrule
\texttt{bg.raised} & \shortstack[l]{0 / \texttt{\#ffffff}\\(1.000, 0.003, 286°)} & \shortstack[l]{950 / \texttt{\#101010}\\(0.173, 0.000, 0°)} & 떠 있는 카드·패널의 바탕 \\
\texttt{bg.canvas} & \shortstack[l]{100 / \texttt{\#f9f9fb}\\(0.982, 0.003, 286°)} & \shortstack[l]{925 / \texttt{\#171717}\\(0.205, 0.000, 0°)} & 페이지와 작업 공간의 바탕 \\
\texttt{bg.sunken} & \shortstack[l]{200 / \texttt{\#eeeef0}\\(0.948, 0.003, 286°)} & \shortstack[l]{900 / \texttt{\#212121}\\(0.248, 0.000, 0°)} & 캔버스보다 한 단계 안쪽인 면 \\
\texttt{bg.well} & \shortstack[l]{400 / \texttt{\#d7d7d9}\\(0.881, 0.003, 286°)} & \shortstack[l]{800 / \texttt{\#2b2b2b}\\(0.289, 0.000, 0°)} & 조작 트랙·깊은 inset 영역 \\
\texttt{border.base} & \shortstack[l]{300 / \texttt{\#e2e2e5}\\(0.915, 0.003, 286°)} & \shortstack[l]{700 / \texttt{\#343434}\\(0.325, 0.000, 0°)} & 기본 경계 \\
\texttt{border.strong} & \shortstack[l]{500 / \texttt{\#ccccce}\\(0.847, 0.003, 286°)} & \shortstack[l]{600 / \texttt{\#444444}\\(0.387, 0.000, 0°)} & 강조 경계 \\
\texttt{border.hover} & \shortstack[l]{600 / \texttt{\#b1b1b3}\\(0.762, 0.003, 286°)} & \shortstack[l]{500 / \texttt{\#5c5c5c}\\(0.475, 0.000, 0°)} & 입력·컨트롤 포인터 강조 \\
\texttt{ink.primary} & \shortstack[l]{950 / \texttt{\#262628}\\(0.271, 0.003, 286°)} & \shortstack[l]{0 / \texttt{\#f5f5f5}\\(0.970, 0.000, 0°)} & 주요 텍스트와 아이콘 \\
\texttt{ink.secondary} & \shortstack[l]{900 / \texttt{\#656567}\\(0.506, 0.003, 286°)} & \shortstack[l]{200 / \texttt{\#adadad}\\(0.748, 0.000, 0°)} & 보조 텍스트 \\
\texttt{ink.tertiary} & \shortstack[l]{700 / \texttt{\#979799}\\(0.677, 0.003, 286°)} & \shortstack[l]{300 / \texttt{\#8e8e8e}\\(0.647, 0.000, 0°)} & 약한 정보와 비활성 표현 \\
\bottomrule
\end{tabularx}
\end{table}

괄호의 OKLCH는 source 입력 좌표이며, HEX는 gamut mapping 후 출력값이다.

위 표에서 step은 두 테마에 동일한 물리적 밝기 의미를 갖는다. Surface token은 그 숫자를 화면 의미에 연결한다. Light에서는 raised가 가장 밝고 well로 갈수록 어두워지며, dark에서는 raised가 가장 어둡고 well로 갈수록 밝아진다. 이는 토큰의 \texttt{raised}가 수치상 더 밝다는 뜻이 아니라, 테마 안에서 콘텐츠와 배경 사이의 표면 역할을 가리키기 때문이다. 앱은 단계 번호 대신 \texttt{bg.raised} 같은 역할 이름을 사용한다.

\subsection{Accent와 interaction}

Accent system은 상태 이름을 palette family와 연결한다. Primary는 neutral에서 만들고, info는 blue, success는 green, warning은 amber, error는 red에 대응한다. 다섯 이름은 동일한 \texttt{solid}, \texttt{wash}, \texttt{soft}, \texttt{subtle} slot을 제공한다. 예를 들어 \texttt{solid}는 의미색 면 또는 강조 잉크에, \texttt{wash}는 solid 위 전경색이나 옅은 배경에 쓰인다\cite{semantics}.

Slot은 색 하나만 고르는 이름이 아니라 색 조합의 계약이다. \texttt{solid}는 채운 면의 배경, \texttt{wash}는 그 면 위 전경이자 아주 옅은 면, \texttt{soft}는 부드러운 면과 hairline 경계, \texttt{subtle}은 muted foreground다. Status accent는 테마에 따라 다른 step을 선택한다. Primary는 neutral의 극성과 surface 체계를 따라 light/dark마다 step을 다르게 고른다. 다음 표는 semantic accent slot의 실제 HEX와 source step을 함께 고정한다.

\begin{table}[htbp]
\centering
\caption{Accent semantic slot의 실제 palette 선택}
\label{tab:accent-semantic-values}
\scriptsize
\begin{tabularx}{\textwidth}{@{}l l *{4}{>{\centering\arraybackslash}X}@{}}
\toprule
Accent & Theme & solid & wash & soft & subtle \\
\midrule
primary & light & \shortstack{\texttt{\#111113}\\950} & \shortstack{\texttt{\#ffffff}\\0} & \shortstack{\texttt{\#d7d7d9}\\300} & \shortstack{\texttt{\#7d7d7f}\\400} \\
primary & dark & \shortstack{\texttt{\#f3f3f5}\\100} & \shortstack{\texttt{\#111113}\\950} & \shortstack{\texttt{\#424244}\\800} & \shortstack{\texttt{\#d7d7d9}\\300} \\
info / blue & light & \shortstack{\texttt{\#2563eb}\\600} & \shortstack{\texttt{\#f5f9fe}\\50} & \shortstack{\texttt{\#d7e5f7}\\150} & \shortstack{\texttt{\#93bbf0}\\300} \\
info / blue & dark & \shortstack{\texttt{\#5694e9}\\400} & \shortstack{\texttt{\#0d255d}\\950} & \shortstack{\texttt{\#102f7a}\\925} & \shortstack{\texttt{\#3a7dea}\\500} \\
success / green & light & \shortstack{\texttt{\#15803d}\\600} & \shortstack{\texttt{\#f5fdf8}\\50} & \shortstack{\texttt{\#d2f0de}\\150} & \shortstack{\texttt{\#86cc9f}\\300} \\
success / green & dark & \shortstack{\texttt{\#35a75f}\\400} & \shortstack{\texttt{\#062b16}\\950} & \shortstack{\texttt{\#09391d}\\925} & \shortstack{\texttt{\#26934e}\\500} \\
warning / amber & light & \shortstack{\texttt{\#b45309}\\600} & \shortstack{\texttt{\#fefbf3}\\50} & \shortstack{\texttt{\#f4e8bb}\\150} & \shortstack{\texttt{\#dbb86f}\\300} \\
warning / amber & dark & \shortstack{\texttt{\#c5861b}\\400} & \shortstack{\texttt{\#421f07}\\950} & \shortstack{\texttt{\#552708}\\925} & \shortstack{\texttt{\#be6d06}\\500} \\
error / red & light & \shortstack{\texttt{\#dc2626}\\600} & \shortstack{\texttt{\#fdf7f7}\\50} & \shortstack{\texttt{\#f6dddd}\\150} & \shortstack{\texttt{\#f4a3a2}\\300} \\
error / red & dark & \shortstack{\texttt{\#ed6969}\\400} & \shortstack{\texttt{\#531313}\\950} & \shortstack{\texttt{\#6b1617}\\925} & \shortstack{\texttt{\#e54c4b}\\500} \\
\bottomrule
\end{tabularx}
\end{table}

Accent step 번호는 neutral step처럼 양 테마에서 밝기 방향이 같지 않다. Palette는 하나의 light-to-dark ramp이고, semantic layer는 테마에 따라 필요한 step을 선택한다. Status accent는 \texttt{solid=600/400}, \texttt{wash=50/950}, \texttt{soft=150/925}, \texttt{subtle=300/500}을 light/dark 순서로 사용한다. Hairline 경계는 별도 slot 없이 \texttt{soft}를 공유한다. Primary는 neutral ramp에서 별도 step을 선택한다.

\subsection{Button, Badge, Alert의 variant별 색 합성}

\path{accentProfile}은 Accent slot을 컴포넌트 variant에 연결한다. 각 profile은 배경, 글자, 경계, hover overlay를 반환한다. CSS emitter는 \texttt{background}, \texttt{color}, \texttt{border-color}를 공통으로 출력하며, interactive class가 있는 경우에만 hover 규칙을 추가한다. Hairline의 border는 hover 중에도 \texttt{soft}를 유지하고, 표면 위에 neutral overlay를 얹는다. Button은 interactive gate를 붙이지만 Badge는 표시용이므로 붙이지 않는다. 따라서 Badge가 pointer 위에 놓여도 hover 색은 발생하지 않는다.

\begin{table}[htbp]
\centering
\caption{Accent variant recipe: (t)는 선택된 accent의 solid·wash·soft·subtle slot}
\label{tab:variant-recipes}
\small
\begin{tabularx}{\textwidth}{@{}l X X X X@{}}
\toprule
Variant & Background & Foreground & Border & Interactive hover \\
\midrule
solid & \texttt{t.solid} & \texttt{t.wash} & transparent & neutral primary면 primarySolid, 나머지는 accentSolid overlay \\
soft & \texttt{t.soft} & \texttt{t.solid} & transparent & neutral overlay \\
hairline & transparent & \texttt{t.solid} & \texttt{t.soft} & neutral overlay \\
ghost & transparent & \texttt{t.solid} & transparent & neutral overlay \\
muted & transparent & \texttt{t.subtle} & transparent & neutral overlay \\
\bottomrule
\end{tabularx}
\end{table}

\texttt{muted}는 accent 매트릭스에서 배경·경계가 투명한 보조 variant다. \texttt{ghost}와 같은 neutral hover를 쓰되 전경에는 선택된 accent의 \texttt{subtle} slot을 사용한다. 따라서 status accent마다 muted 글자에도 해당 색조가 유지된다. 반면 \texttt{solid}의 hover는 기존 바탕색을 교체하지 않는다. CSS가 overlay를 단색 linear-gradient로 얹어 원래 pigment를 유지하면서 밝기만 조금 낮춘다. 이 방식은 빨강·초록·파랑 같은 accent에 검정 HEX를 직접 더하지 않게 한다.

같은 recipe의 결과를 실제 컴포넌트가 어떻게 쓰는지는 다음과 같다. Button의 default는 primary/solid이고 accent와 variant prop으로 resolver를 고른다. Badge는 primary/soft를 기본으로 하며 상태 표시에는 info, success, warning, error를 주로 사용한다. Alert는 info/soft를 기본으로 하고, 숨겨진 접두어 텍스트도 함께 내보내 색 이외의 상태 단서를 제공한다. Banner는 primary/solid 또는 상태별 solid 면을 만들고 반드시 그 accent의 wash 글자를 쓴다. Link는 info/solid를 사용한다. 이 구성은 토큰 선택과 geometry를 분리한다. 버튼의 반경과 높이는 Button CSS가 정하고 색은 accent-variant class profile이 정한다.

\subsection{Form control과 본문 텍스트의 색 상태}

Field, Select, Textarea는 \path{bg.raised} 면에 \path{border.base} 경계를 두고 \path{ink.primary}로 입력값을 표시한다. Placeholder는 \path{ink.tertiary}, Field adornment는 \path{ink.secondary}를 사용한다. 기본 pointer hover는 \path{border.hover}로 경계를 또렷하게 한다. 클릭한 뒤의 pointer focus는 \path{ink.secondary} 경계로 표현한다. 키보드 \path{:focus-visible}에는 상태색 \path{interaction.focus}를 외곽 ring으로 적용하고 입력 자체의 경계와 구분한다.

Validation status가 지정되면 해당 status accent의 \texttt{solid}가 border가 되며, hover에는 같은 색 family의 반투명 glow가 붙는다. Validation 상태가 키보드 focus와 동시에 켜지면 glow를 끄고 focus ring만 남긴다.

Checkbox와 Radio는 꺼졌을 때 \texttt{bg.raised}와 \texttt{border.strong}을 쓴다. Checked Checkbox는 primary solid 면 위에 primary wash 체크 표시를 그린다. Radio는 선택될 때 edge를 primary solid로 바꾸고 중심점을 같은 색으로 채운다. Switch는 꺼졌을 때 \texttt{bg.well} 트랙과 \texttt{bg.raised} 손잡이, 켜졌을 때 primary solid 트랙을 사용한다. 세 컨트롤의 keyboard focus는 모두 같은 info 계열 \texttt{interaction.focus} ring을 쓴다. Disabled opacity는 색 토큰을 다시 만들지 않고 전체 요소 opacity로 낮춘다.

Text의 기본 variant 색은 \path{TEXT_CONTENT} mapping이 고른다. Display, Title, Heading, Subheading, Body는 \path{ink.primary}를, Caption과 Label은 \path{ink.secondary}를 사용한다. 별도 ink prop은 neutral 단계 또는 상태 accent solid를 고를 수 있고 inherit는 부모가 정한 foreground를 그대로 따른다. 이 때문에 Banner나 Alert 내부 텍스트는 자기 기본 잉크로 덮어쓰지 않고 면에 맞는 wash 전경색을 상속한다.

Interaction system은 중립 hover, primary solid hover, accent solid hover, 상태별 glow, selected, active, focus, scrim을 다룬다. Focus는 상호작용 상태이므로 accent palette를 직접 고르지 않고 \texttt{interaction.focus}로 공개하며 blue 600 palette token을 선택한다. Focus ring의 폭과 offset은 색이 아닌 geometry이므로 shape system에 남는다. 이 분리는 색의 역할과 형태 수치를 섞지 않는 계약이다\cite{architecture,semantics}.

\begin{longtable}{@{}l p{0.33\textwidth} p{0.33\textwidth}@{}}
\caption{Interaction 및 shadow 색 source와 실제 8-digit HEX}
\label{tab:interaction-values}\\
\toprule
Token & Light: source / opacity / output & Dark: source / opacity / output \\
\midrule
\endfirsthead
\multicolumn{3}{l}{\tablename\ \thetable\ (continued)}\\
\toprule
Token & Light: source / opacity / output & Dark: source / opacity / output \\
\midrule
\endhead
\midrule
\multicolumn{3}{r}{Continued on next page}\\
\endfoot
\bottomrule
\endlastfoot
hover.neutral & neutral.950 / 0.05 / \texttt{\#2626280d} & neutral.0 / 0.036 / \texttt{\#f5f5f509} \\
hover.primarySolid & neutral.0 / 0.12 / \texttt{\#ffffff1f} & neutral.950 / 0.12 / \texttt{\#1010101f} \\
hover.accentSolid & neutral.950 / 0.14 / \texttt{\#26262824} & neutral.950 / 0.14 / \texttt{\#10101024} \\
selected & neutral.950 / 0.13 / \texttt{\#26262821} & neutral.0 / 0.10 / \texttt{\#f5f5f51a} \\
active & neutral.950 / 0.16 / \texttt{\#26262829} & neutral.0 / 0.13 / \texttt{\#f5f5f521} \\
focus & blue.600 / -- / \texttt{\#2563eb} & blue.600 / -- / \texttt{\#5492ec} \\
scrim & neutral.950 / 0.32 / \texttt{\#26262852} & neutral.950 / 0.58 / \texttt{\#10101094} \\
hover.statusGlow.info & blue.600 / 0.22 / \texttt{\#2563eb38} & blue.600 / 0.22 / \texttt{\#5492ec38} \\
hover.statusGlow.success & green.600 / 0.22 / \texttt{\#15803d38} & green.600 / 0.22 / \texttt{\#32ac6438} \\
hover.statusGlow.warning & amber.600 / 0.22 / \texttt{\#b4530938} & amber.600 / 0.22 / \texttt{\#c5830038} \\
hover.statusGlow.error & red.600 / 0.22 / \texttt{\#dc262638} & red.600 / 0.22 / \texttt{\#df676838} \\
shadow.overlayStrong & neutral.950 / 0.075 / \texttt{\#26262813} & neutral.950 / 0.62 / \texttt{\#1010109e} \\
shadow.overlaySoft & neutral.950 / 0.035 / \texttt{\#26262809} & neutral.950 / 0.50 / \texttt{\#10101080} \\
shadow.overlayMinimal & neutral.950 / 0.08 / \texttt{\#26262814} & neutral.950 / 0.45 / \texttt{\#10101073} \\
\end{longtable}

Hover와 selected overlay는 표면색을 token 값으로 교체하지 않는다. \path{color.css.ts}가 pseudo-class에서 같은 8-digit HEX를 \texttt{linear-gradient(color, color)}로 얹어 기존 바탕을 보존한다. Text 버튼과 solid status 버튼은 interactive 상태에서만 hover 규칙을 받는다. Scrim은 별도 전체화면 overlay 색이고, status glow는 \texttt{Field}의 validation border 주위에만 shape의 12px blur를 빌려 사용한다. Alpha 값은 light와 dark의 대비를 각각 검토해 정하므로 한 모드의 수치를 반대 모드에 기계적으로 복사하지 않는다.

Constants system은 shadow tint처럼 palette에서 생성되지만 interaction 자체는 아닌 고정 효과 색을 보관한다. 그림자의 거리, blur, spread는 shape/effect geometry에 속하고, 그림자에 입힐 tint만 color constants에 둔다. 모든 고정 효과값을 무조건 constants에 넣는 것이 아니라, 색의 출처와 사용 역할이 semantic branch에 들어맞는지 기준으로 분류한다.

Shadow 최종값은 \path{theme.css.ts}가 shape geometry와 color constant를 조립한다. Light overlay는 \texttt{0 8px 24px 0}의 강한 layer와 \texttt{0 2px 6px 0}의 부드러운 layer를 쌓고, dark는 \texttt{0 8px 22px -5px}와 \texttt{0 3px 7px -2px}를 쓴다. Switch thumb의 minimal shadow도 light \texttt{0 1px 3px 0}, dark \texttt{0 1px 2px 0}으로 기하를 달리한다. 각 layer의 tint는 바로 앞 표의 semantic color constant를 선택하므로 밝은 쪽의 검정 shadow와 어두운 쪽의 더 짙은 검정 shadow를 별도 CSS HEX로 컴포넌트에 박지 않는다. Focus ring은 info solid을 색으로, 2px 폭과 2px offset을 geometry로 나누어 가진다.

\section{CSS 노출과 실제 사용}

색이 브라우저에 도달하는 과정은 세 단계다.

\begin{enumerate}
\item \texttt{palette.ts}의 OKLCH 입력을 \texttt{color-ladder.ts}가 palette HEX tree로 생성한다. 이 단계는 private source/generator 경계이며 semantic 역할은 모른다.
\item \texttt{colors.ts}가 light/dark의 semantic tree를 만들고 모든 leaf를 생성된 palette token으로 선택한다. \texttt{tokens/index.ts}의 \texttt{buildVarTree}는 공개 token API를 변수 참조 형태로 만든다. 예를 들어 \path{tokens.color.bg.canvas}는 HEX 값이 아니라 \path{var(--pui-color-bg-canvas)}를 가리킨다.
\item \path{theme.css.ts}가 semantic tree를 \path{walkValues}로 순회하고 경로를 \path{pathToCssVar}로 CSS 이름에 변환한다. Light는 \path{:root} 기본값으로도 제공하고, 각 테마 scope가 값을 교체한다. 같은 변수 이름의 값만 바뀌므로 component CSS는 mode 분기를 작성하지 않는다.
\end{enumerate}

예를 들어 \path{tokens.color.bg.canvas}는 CSS custom property를 통해 light에서 neutral.100 (\texttt{\#f9f9fb}), dark에서 neutral.925 (\texttt{\#171717})를 참조한다. \path{utility.css.ts}의 surface utility도 같은 네 bg token을 사용한다. 앱은 임의의 palette step을 선택하지 않는다. 테마 전환은 palette 재생성이나 컴포넌트 색 계산이 아니라 CSS scope 변경이다\cite{semantics,theme,tokenindex,utility}.

Accent component는 별도 dispatch 단계를 거친다. \path{resolveAccent(accent, variant)}는 accent와 variant 조합에 대응하는 class를 반환한다. \path{color.css.ts}는 ACCENTS x VARIANTS 조합을 전부 순회해 CSS를 생성하고, 각 조합의 수식은 \path{accentProfile}이 accent slot과 interaction variable을 조합해 제공한다. 그래서 런타임에는 class 이름으로 이미 빌드된 CSS를 선택하고 color math나 HEX lookup은 실행하지 않는다. Hover selector는 \path{.pui-interactive}가 함께 있을 때만 활성화된다. Button에는 gate가 있고 Badge에는 없어 같은 accent/variant라도 hover 유무가 분리된다\cite{profile,emitter,resolver}.

앱 사용 계약도 같은 방향이다. Doodle, daily-todo, booklog는 Paper UI package의 token/atom을 소비하므로 공용 semantic mapping 변경을 함께 받는다. 다만 보고서의 검증 대상은 현재 저장소의 Paper UI 구현이다. 앱이 package를 실제 import하는 경로와 각 앱의 로컬 CSS override까지 이 문서가 전수 검사한다는 뜻은 아니다. palette 예시를 보여 주는 Palette Playground는 source 값을 검사하는 의도된 예외다.

알려진 우회값도 기록할 필요가 있다. 현재 \texttt{Select.css.ts}의 화살표는 CSS data URI 내부 SVG stroke를 \texttt{\%2371717a}, \texttt{\%23a1a1aa}, \texttt{\%238b8b92}로 직접 고정한다. 이 값은 semantic variable이 아니고 dark/placeholder selector도 별도 값으로 바꾸므로 tokens-only 계약을 완전히 따르지 않는 구현 예외다. data URI 안에서 CSS variable을 직접 주입하기 어렵다는 구현상 이유가 있지만, inline SVG/mask/currentColor 기반 아이콘이나 별도 토큰 구동 자산으로 옮겨야 할 debt다. 이 보고서는 해당 코드를 고치지 않고 구현의 사실로 공개한다.

이 구조에서 팔레트 변경은 실제 제품 화면에 도달할 때까지 같은 경로를 통과한다. 최근 라이트 neutral 조정은 palette source의 hue/chroma만 바꾸고, semantic mapping이나 Doodle·daily-todo·booklog의 개별 색상값은 추가하지 않았다. 따라서 세 앱의 surface와 ink는 모두 같은 배포 패키지에서 갱신된다. 앱별로 HEX를 덮어쓰면 토큰 시스템의 진실 공급원이 둘로 갈라지므로, 제품별 차이가 필요할 때는 새 semantic 역할이나 공식 theme profile이 필요한지 먼저 검토해야 한다.

\section{정합성 계약과 검증}

현재 자동 검사는 구조적 계약을 중심으로 한다. \texttt{colors.test.ts}는 light/dark 양쪽에서 neutral 13단계와 accent 11단계가 완전한지, neutral 0이 white이고 950이 near-black인지, light neutral의 lightness가 단조롭게 진행하는지, 출력이 6자리 또는 8자리 HEX인지 확인한다. 이어 모든 semantic leaf가 실제 생성된 palette token과 같고, surface·border·ink가 문서화된 palette step을 선택하는지 검사한다. Accent의 status-to-palette mapping과 각 slot step, focus source, alpha를 포함한 interaction 및 shadow token도 검사한다\cite{colortests}.

\texttt{color-profile.test.ts}는 muted가 모든 accent에서 같은 Ghost형 투명 면, secondary ink, neutral hover를 쓰는지 확인한다. CSS emitter는 resolver의 여섯 variant와 다섯 accent 조합을 생성한다. 이는 계산된 CSS profile의 계약이며 브라우저 screenshot이나 실제 명암비 검사는 아니다\cite{profiletests}.

Lint와 CSS architecture는 잘못된 접근 경로를 막는다. Palette·ladder·semantic value를 소비 레이어에서 직접 import하지 못하게 하고, raw 색값은 token source 이외에 작성하지 못하게 한다. Component와 app은 semantic CSS variable을 사용한다. Palette Playground는 원천 좌표와 생성 결과를 함께 보여주는 의도된 inspector 예외다\cite{architecture,tokensdoc}.

이 검사는 색의 계층과 참조 정합성을 보장하지만 모든 전경·배경 조합의 명도 대비 비율을 평가하지는 않는다. 특히 tertiary ink, disabled opacity가 적용된 컨트롤, 아주 작은 아이콘과 텍스트, translucent overlay가 겹치는 상황은 실제 합성 결과를 별도로 검사해야 한다. 또한 Select SVG처럼 lint가 raw color를 차단하더라도 data URI 문자열 내부를 놓칠 수 있는 예외가 남아 있다. 다음 단계로 semantic role pair의 contrast profile과 예외 목록을 정의하고, 라이트·다크 양 테마에서 최소 대비 기준 및 우회값을 자동 탐지하는 검사를 추가하는 편이 타당하다. 검사는 토큰을 무조건 밝히거나 어둡게 만들기보다 기준을 벗어난 조합을 구체적으로 드러내야 한다.

\section{결론}

Paper UI의 Color System은 팔레트 보관소보다 색을 통제하는 경계에 가깝다. OKLCH source는 hue·chroma·lightness의 정체성을 소유하고, HEX generator는 브라우저에 배포할 palette token을 만든다. Semantic layer는 그 값에 UI 역할을 부여하고, theme CSS와 컴포넌트는 역할 이름만 소비한다. 이 구조는 팔레트 변화가 여러 앱에 일관되게 전달되면서도, 컴포넌트가 원천 색을 임의로 선택하는 일을 줄인다.

현재 구현의 강점은 생성기, semantic mapping, CSS 소비 경계가 분리되어 있어 source 좌표부터 실제 variant의 합성 결과까지 추적 가능하다는 점이다. 한계도 같은 구체성으로 드러난다. CSS data URI의 직접 stroke 값은 계약의 예외이며, 구조 테스트만으로는 화면의 contrast 품질을 증명하지 못한다. 따라서 색 시스템을 완전히 이해한다는 것은 token table만 아는 것이 아니라, 이 변환 경로와 variant recipe, 런타임 scope, 검증의 범위 및 명시적 예외를 함께 아는 것이다.

\begin{thebibliography}{9}
\bibitem{architecture} Studio Baeks, ``Color architecture,'' \textit{Paper UI repository}, \path{docs/foundations/color-architecture.md}, 2026.
\bibitem{palette} Studio Baeks, ``OKLCH palette source,'' \textit{Paper UI repository}, \path{workspace/packages/paper-ui/src/tokens/palette.ts}, 2026-09-25 snapshot.
\bibitem{ladder} Studio Baeks, ``OKLCH to HEX palette-token generator,'' \textit{Paper UI repository}, \path{workspace/packages/paper-ui/src/tokens/color-ladder.ts}, 2026-09-25 snapshot.
\bibitem{semantics} Studio Baeks, ``Semantic color systems,'' \textit{Paper UI repository}, \path{workspace/packages/paper-ui/src/tokens/colors.ts}, 2026-09-25 snapshot.
\bibitem{theme} Studio Baeks, ``Theme CSS generation,'' \textit{Paper UI repository}, \path{workspace/packages/paper-ui/src/styles/theme.css.ts}, 2026-09-25 snapshot.
\bibitem{colortests} Studio Baeks, ``Color-system contract tests,'' \textit{Paper UI repository}, \path{workspace/packages/paper-ui/src/tokens/colors.test.ts}, 2026-09-25 snapshot.
\bibitem{profile} Studio Baeks, ``Accent variant recipes,'' \textit{Paper UI repository}, \path{workspace/packages/paper-ui/src/styles/color-profile.ts}, 2026-09-25 snapshot.
\bibitem{emitter} Studio Baeks, ``Accent CSS emitter,'' \textit{Paper UI repository}, \path{workspace/packages/paper-ui/src/styles/color.css.ts}, 2026-09-25 snapshot.
\bibitem{resolver} Studio Baeks, ``Color class resolver,'' \textit{Paper UI repository}, \path{workspace/packages/paper-ui/src/resolvers/color.ts}, 2026-09-25 snapshot.
\bibitem{profiletests} Studio Baeks, ``Accent profile tests,'' \textit{Paper UI repository}, \path{workspace/packages/paper-ui/src/styles/color-profile.test.ts}, 2026-09-25 snapshot.
\bibitem{tokenindex} Studio Baeks, ``Public token variable tree,'' \textit{Paper UI repository}, \path{workspace/packages/paper-ui/src/tokens/index.ts}, 2026-09-25 snapshot.
\bibitem{utility} Studio Baeks, ``Surface utility styles,'' \textit{Paper UI repository}, \path{workspace/packages/paper-ui/src/styles/utility.css.ts}, 2026-09-25 snapshot.
\bibitem{tokensdoc} Studio Baeks, ``Token architecture,'' \textit{Paper UI repository}, \path{docs/foundations/tokens.md}, 2026.
\end{thebibliography}

\end{document}
